% H3 Hiemenz — corrected question statement, XeLaTeX + xepersian
\documentclass[12pt,a4paper]{article}
\usepackage{amsmath,amssymb}
\usepackage{booktabs}
\usepackage{graphicx}
\usepackage{geometry}
\geometry{margin=25mm}
\usepackage[hidelinks]{hyperref}
\usepackage{xepersian}
\settextfont{Vazirmatn}
\setdigitfont{Vazirmatn}
\usepackage{setspace}
\setstretch{1.25}
\begin{document}

\begin{center}
{\Large\bfseries صورت سؤال اصلاح‌شده\par}
\vspace{1ex}
{\large جریان نقطهٔ سکون هایمنز؛ حل تشابهی لایهٔ مرزی\par}
\vspace{1ex}
{\small درس لایهٔ مرزی · نسخهٔ اصلاح‌شدهٔ یادداشت دستی\par}
\end{center}

\noindent\textbf{یادداشت اصلاح:} این متن، نسخهٔ تمیز و اصلاح‌شدهٔ یادداشت
دستی است. تنها اصلاحِ محتوایی، علامت فرمول‌های سرعت است (بخش ۴)؛ بقیهٔ
فرمول‌ها همان یادداشت‌اند. در یادداشت، لزجت سینماتیکی با \lr{$\gamma$}
نوشته شده است؛ از نظر بُعدی همان \lr{$\nu$} است.

\section{مسئله}
جریان پایا، دوبعدی و تراکم‌ناپذیرِ یک سیال لزج را روی یک صفحهٔ تخت در
نظر بگیرید. دیوار روی \lr{$x_2=0$} است و جریان بیرونیِ نقطهٔ سکون،
\lr{$U_e = a x_1$} با \lr{$a>0$} است؛ یعنی سرعت بیرونی با فاصله از نقطهٔ
سکون خطی زیاد می‌شود. \lr{$x_1$} راستای مماس بر دیوار، \lr{$x_2$} راستای
عمود بر دیوار، و \lr{$u_1$} و \lr{$u_2$} مؤلفه‌های سرعت در همین دو راستا
هستند. لزجت سینماتیکی را \lr{$\gamma$} بگیرید.

\section{متغیرهای بی‌بعد}
طول مشخصه \lr{$\sqrt{\gamma/a}$} و سرعت مشخصه \lr{$\sqrt{\gamma a}$}
هستند. متغیرهای بی‌بعد (با ستاره) چنین تعریف می‌شوند:
\begin{equation}
x_1^*, x_2^* = \frac{x_1, x_2}{\sqrt{\gamma/a}},
\qquad
u_1^*, u_2^* = \frac{u_1, u_2}{\sqrt{\gamma a}}.
\end{equation}
متغیر تشابهی، همان مختصهٔ بی‌بعدِ عمودی است:
\lr{$\eta = x_2^*$}.

\section{فرض تشابهی (قرارداد استاندارد)}
فرض کنید سرعت‌ها فقط از راه یک تابعِ تک‌متغیرهٔ \lr{$F(\eta)$} بیان شوند:
\begin{equation}
u_1^* = x_1^*\, F'(\eta),
\qquad
u_2^* = -F(\eta),
\qquad \eta = x_2^*.
\end{equation}
در اینجا \lr{$F'$} مشتق نسبت به \lr{$\eta$} است.

\section{خواسته‌ها}
\begin{enumerate}
\item نشان دهید فرض بالا، معادلهٔ پیوستگیِ بی‌بعد،
\lr{$\partial u_1^*/\partial x_1^* + \partial u_2^*/\partial x_2^* = 0$}،
را برقرار می‌کند.
\item با جای‌گذاری فرض بالا در معادلهٔ مومنتمِ لایهٔ مرزیِ بی‌بعد،
\lr{$u_1^* \partial u_1^*/\partial x_1^* + u_2^* \partial u_1^*/\partial x_2^*
= x_1^* + \partial^2 u_1^*/\partial {x_2^*}^2$}،
معادلهٔ زیر را به دست آورید:
\begin{equation}
F''' + F F'' - (F')^2 + 1 = 0.
\end{equation}
\item شرایط مرزی زیر را بنویسید و هر یک را تفسیر کنید:
\begin{equation}
F(0) = 0, \qquad F'(0) = 0, \qquad F'(\infty) = 1.
\end{equation}
کدام شرط، چسبندگی و کدام‌یک، نفوذناپذیری دیوار است؟ شرط دوردست،
چه می‌گوید؟
\item معادلهٔ بالا را با شرایط مرزی داده‌شده به‌صورت عددی حل کنید و
مقدار \lr{$F''(0)$} را گزارش کنید. این مقدار با چه کمیت فیزیکی‌ای
(تنش برشی دیوار) متناسب است؟
\end{enumerate}

\section{اصلاح نسبت به یادداشت دستی}
در یادداشت دستی، فرمول‌های سرعت با علامت برعکس آمده بودند:
\lr{$u_2^* = F$} و \lr{$u_1^* = -x_1^* F'$}. آن دو با معادله و شرط
دوردستِ بالا سازگار نیستند؛ با آن علامت‌ها، معادله
\lr{$F''' - F F'' + (F')^2 - 1 = 0$} و شرط \lr{$F'(\infty) = -1$}
به دست می‌آید. در این نسخه، قرارداد استانداردِ هایمنز گرفته شده است
تا جواب با جدول‌های مرجع بخواند.

\end{document}
